\documentclass[10pt,a4paper]{article}

\usepackage[utf8]{inputenc}
\usepackage[T1]{fontenc}
\usepackage{amsmath,amssymb,amsthm}
\usepackage{booktabs}
\usepackage{tabularx}
\usepackage{enumitem}
\usepackage{hyperref}
\usepackage{natbib}
\usepackage[margin=2cm]{geometry}

\newtheorem{lemma}{Lemma}

\title{Where Counterfactual Invariance Helps, and Where It Doesn't:\\
A Three-Axis Evaluation of Antisemitism Detection on Text}
\author{Giacomo Case \and Yasmin Chaqor \and Leonardo Cartesegna\\
EPFL --- EE-559 Deep Learning Mini-Project}
\date{Spring 2026}

\begin{document}

\maketitle

\begin{abstract}
We study antisemitism detection in English tweets along three connected axes
--- positive-class F\textsubscript{1}, group-conditional Expected Calibration
Error (ECE), and counterfactual stability under identity-token swaps within a
religiously-symmetric vocabulary. We propose \textbf{Counterfactual Calibration
Invariance (CCI v2)}, a training-time loss that penalises Jensen-Shannon
divergence between a tweet's prediction and the prediction on its identity-swapped
counterfactual, modulated by an optional per-group temperature head. Across three
architectures (DeBERTa-v3-base, DeBERTa-v3-large, RoBERTa-large; 6 seeds each)
CCI v2 reduces the hard counterfactual flip rate (CFR\textsubscript{hard}) by
49--73~\% versus a strong focal+identity-mask baseline with positive-class
F\textsubscript{1} preserved or improved (best 0.72). A short structural Lemma
shows that scalar Temperature Scaling cannot reduce CFR\textsubscript{hard} at
the decision threshold $1/2$; we verify this bit-exactly across six seeds.
A separate finding: applying multicalibration~\citep{hebertjohnson2018multicalibration}
post-hoc cuts CFR\textsubscript{hard} by 8~\% but \emph{inflates the
L\textsubscript{1}-norm soft flip rate by 96~\%} --- a tension between two
well-respected fairness desiderata. Prompted LLMs (Llama-3.1-8B, Mistral-7B,
Qwen-2.5-7B) under zero-shot, few-shot, and Guided
Chain-of-Thought~\citep{patel2025evaluating} reach F\textsubscript{1}\textsuperscript{+}
$\le 0.445$ --- 22 points below our 184M-parameter fine-tuned model.
Cross-dataset transfer to the Jewish-target subsets of HateXplain and
ToxiGen exposes a clean threshold-shift artefact: at the default decision
threshold $t{=}0.5$, CCI v2's score distribution sits at a systematically
more conservative operating point (median optimal $t^*$ between 0.05 and 0.15)
and underperforms on 4 of 6 architecture/dataset cells; at each model's
PR-curve-optimal threshold CCI v2 wins or ties on 5 of 6 cells (the lone
loss is $-0.011$, within seed noise) and PR-AUC is preserved or improved
on every iso-architecture cell (per-cell $\Delta$PR-AUC range $[-0.007, +0.021]$).
A source-tuned threshold check ($n{=}68$ cells; tune on
GoldStandard2024's test split, apply unchanged to OOD) locates the residual
gap in a $\approx 0.4$-unit prevalence-shift on the optimal threshold,
not in the representation: the in-distribution counterfactual-invariance
gain transfers as preserved ranking quality, and OOD-aware threshold
calibration converts that ranking into deployable decisions.
\end{abstract}

\section{Research Questions}

Rather than ask a single composite question, we decompose the project into
five concrete sub-questions, each answered in a specific results section.
The decomposition lets us cleanly attribute positive and negative findings
to specific scientific claims rather than to one method.

\begin{description}[leftmargin=*,itemsep=4pt]
    \item[\textbf{RQ1 (are these three qualities really separate? \S\ref{sec:results-id})}]
    Imagine a classifier that flags a tweet as antisemitic. We can ask three
    different things about it: (i) how often is it \emph{right}? (ii) when it
    reports ``70~\% confident,'' does that actually happen 70~\% of the time
    --- and is that still true when we look only at tweets about a particular
    identity group? (iii) if we rewrite the same tweet swapping ``Jew'' for
    ``Muslim'' or ``Christian,'' does the verdict stay the same? Are these
    three properties \emph{independent dials} we have to tune separately,
    or is one underlying ``quality'' driving all three so that pushing
    accuracy upward gives us the others for free? \emph{Technical
    restatement:} are positive-class F\textsubscript{1}, group-conditional
    Expected Calibration Error, and counterfactual flip rate under
    identity-token swaps three orthogonal axes of classifier behaviour, or
    a single latent factor?

    \item[\textbf{RQ2 (fix it during training, or patch it afterwards? \S\ref{sec:lemma}--\ref{sec:tension})}]
    Suppose a classifier flips its verdict on a tweet when we swap one
    identity word for another. To stop that, do we have to retrain the model
    (the slow, expensive option), or can we leave the trained model alone
    and just rescale its confidence scores afterwards (a much cheaper fix)?
    Two well-known after-the-fact fixes exist: \emph{Temperature Scaling},
    which divides every confidence score by a single learned constant, and
    \emph{multicalibration}~\citep{hebertjohnson2018multicalibration}, which
    adjusts scores separately for each identity subgroup. Is either enough,
    or does the fix have to be baked into training? \emph{Technical
    restatement:} for reducing hard counterfactual flip rate at the
    decision threshold $\tfrac{1}{2}$, is post-hoc calibration (Temperature
    Scaling, multicalibration) sufficient, or is a training-time loss
    required?

    \item[\textbf{RQ3 (does the fix carry over to bigger and different models? \S\ref{sec:results-id})}]
    If a training trick works on one neural network, does it still work
    when we use a network that is twice as big but built the same way
    (think: the same car model in a larger engine size)? And does it work
    on a network from a different manufacturer entirely (different
    word-splitting rules, different internal architecture)? \emph{Technical
    restatement:} do counterfactual-invariance gains measured on
    DeBERTa-v3-base transfer to DeBERTa-v3-large (same family, $2.4\times$
    parameters) and to RoBERTa-large (different family: different tokenizer,
    different attention implementation)?

    \item[\textbf{RQ4 (does the fix carry over to unfamiliar data? \S\ref{sec:results-ood})}]
    We train and evaluate the fix on one dataset. When we then test it on a
    \emph{different} dataset --- collected by different researchers,
    labelled by different rules, with a different mix of tweets --- does
    the improvement still show up? More specifically: our fix teaches the
    model to be insensitive to a fixed list of identity words. Does that
    insensitivity generalise to new data, or only to data that happens to
    contain words from that list? \emph{Technical restatement:} under what
    conditions do in-distribution counterfactual-invariance gains transfer
    out-of-distribution (HateXplain-Jewish, ToxiGen-Jewish), and how does
    transfer depend on the lexical overlap between the target corpus and
    the swap vocabulary?

    \item[\textbf{RQ5 (can a large chatbot do this job out of the box? \S\ref{sec:results-llm})}]
    Big instruction-following language models like Llama or Mistral can be
    asked, in plain English, ``is this tweet antisemitic?'' --- with no
    training, with a few worked examples, or with a step-by-step reasoning
    prompt. Can such an off-the-shelf model match a much smaller specialist
    classifier that we fine-tuned specifically for this task? \emph{Technical
    restatement:} can prompted 7--8B-parameter instruction-tuned LLMs
    (Llama-3.1-8B, Mistral-7B, Qwen-2.5-7B) under zero-shot, few-shot, or
    Guided Chain-of-Thought~\citep{patel2025evaluating} prompting match or
    exceed our 184M-parameter fine-tuned DeBERTa-v3-base baseline?
\end{description}

\textbf{Anticipated answers (in plain English):}

\begin{itemize}[leftmargin=*,itemsep=3pt]
\item \textbf{RQ1:} \emph{The three properties really are separate.} We give
    a short mathematical proof that the cheapest post-hoc fix (Temperature
    Scaling) cannot reduce the rate at which a classifier flips its verdict
    on swapped tweets. Going the other direction, a more sophisticated
    post-hoc fix (multicalibration) does improve group-by-group calibration
    but \emph{worsens} the average confidence shift on swapped tweets by
    96~\%. So no single optimisation target captures all three.
\item \textbf{RQ2:} \emph{Patching afterwards is not enough --- you have to
    retrain.} Multicalibration cuts the hard flip rate by only about 8~\%,
    and at the cost of the soft flip rate (RQ1). Temperature Scaling
    provably cannot move it at all (Lemma~\ref{lem:ts-cfr}).
\item \textbf{RQ3:} \emph{Yes, on the original dataset, the fix carries
    over.} Our proposed training-time loss, Counterfactual Calibration
    Invariance (CCI v2), reduces the swap-induced flip rate by
    49\,/\,55\,/\,73~\% on DeBERTa-v3-base (184M parameters),
    DeBERTa-v3-large (440M, same family), and RoBERTa-large (355M,
    different family) respectively.
\item \textbf{RQ4:} \emph{The cross-dataset gap is a decision-threshold
    artefact, not a representation gap.} At the default cutoff
    $t{=}0.5$, CCI v2 loses positive-class F\textsubscript{1} on 4 of
    6 architecture/dataset cells because its training compresses scores
    toward the centre of the swap-pair, leaving it at a more conservative
    operating point cross-dataset. Yet PR-AUC (a threshold-free measure
    of ranking quality) is preserved or improved on every iso-architecture
    cell (per-cell $\Delta$ range $[-0.007, +0.021]$, CCI v2 winning on
    4 of 6). When each model is evaluated at its own PR-curve-optimal
    threshold (typically $t^* \in [0.05, 0.18]$), CCI v2 wins or ties on
    5 of 6 cells; the one loss is $-0.011$, well within seed noise.
    Out-of-distribution \emph{ranking quality} therefore does transfer;
    out-of-distribution \emph{decision quality at a fixed default threshold}
    does not, and the gap is closed by threshold tuning rather than by
    a different training objective.
\item \textbf{RQ5:} \emph{No.} The best-performing prompted large model
    (Llama-3.1-8B with Guided Chain-of-Thought) reaches a positive-class
    F\textsubscript{1} of 0.445; the fine-tuned baselines reach 0.660--0.719
    --- a gap of 22--27 points, replicated across three different LLM
    families.
\end{itemize}

\section{Method}

\subsection{Dataset and three-axis evaluation}

We use the GoldStandard2024 dataset~\citep{jikeli2024goldstandard}: 11{,}311
IHRA-annotated~\citep{ihra2016working} English tweets, of which 17.3~\% are
labelled antisemitic, with four keyword sub-populations (Jews, Israel, Kikes,
ZioNazi) at very different base rates (11.4~\%, 16.1~\%, 34.3~\%, 88.3~\%).
After deduplication the test split contains 1{,}060 tweets (187 positives).
We evaluate models along three axes:

\begin{itemize}[leftmargin=*,itemsep=1pt]
    \item \textbf{Accuracy:} positive-class F\textsubscript{1} (F\textsubscript{1}\textsuperscript{+}).
    \item \textbf{Calibration:} group-conditional ECE with 15
    equal-width bins~\citep{kumar2019verified}, plus worst-group ECE.
    \item \textbf{Counterfactual stability:} CFR\textsubscript{hard} =
    $\Pr[\mathbb{1}\{p(x){\geq}\tfrac{1}{2}\}\neq\mathbb{1}\{p(x'){\geq}\tfrac{1}{2}\}]$,
    plus its soft variants CFR\textsubscript{soft, L\textsubscript{1}}
    $= \mathbb{E}\,|p(x){-}p(x')|$ and the binary Jensen-Shannon
    counterpart CFR\textsubscript{soft, JS}. Pairs $(x, x')$ are produced
    by an algorithmic swap engine that substitutes religious-ethnic Jewish
    tokens (\emph{jew}, \emph{jews}, \emph{jewish}, \emph{hebrew},
    \emph{judaism}) with one of five alternative groups (\emph{Muslim},
    \emph{Christian}, \emph{Hindu}, \emph{Buddhist}, \emph{Atheist}); a
    heuristic symmetry filter rejects swaps involving slurs or geopolitical
    proxies. The fixed test-time manifest contains 2{,}117 pairs.
\end{itemize}

\subsection{CCI v2 loss}

For input $x_i$ with label $y_i$ and counterfactual partners $(x'_{i,k})_{k=1}^{K_i}$,
let $z(x) = \mathrm{logit}_1(x) - \mathrm{logit}_0(x)$. We minimise
\[
\mathcal{L}(x_i)=\text{Focal}(f(x_i), y_i)
+\lambda\cdot\frac{1}{K_i}\sum_{k=1}^{K_i}
\mathrm{JS}_b\!\left(\sigma\!\bigl(z(x_i)/T_{g_i}\bigr),
                     \sigma\!\bigl(z(x'_{i,k})/T_{g_i}\bigr)\right),
\]
where $\sigma$ is the sigmoid, $\mathrm{JS}_b$ is the binary Jensen-Shannon
divergence, $T_{g_i}$ is the temperature for example $i$'s identity group
(one of 5), and $\lambda{=}1$. CCI v2 generalises Counterfactual Logit
Pairing~\citep{garg2019counterfactual,davani2021improving} along two axes:
divergence in \emph{probability} space (rather than $L_p$ in logit space)
and a per-group temperature head. Our headline (\emph{pivot}) cell fixes
$T_g{=}1$ everywhere; the learnable-$T$ variant is kept as an ablation
comparator after an in-cell comparison showed it does not improve
CFR\textsubscript{hard} but loses on FNED.

\subsection{Baselines and architectures}

\textbf{Model C} (strong baseline): DeBERTa-v3-base + focal loss
($\gamma{=}2.0,\,\alpha{=}0.75$) + identity-token masking augmentation
($p{=}0.3$). \textbf{Davani CLP}~\citep{davani2021improving}: same backbone
+ $L_1$ pair-loss in logit space at $\lambda\in\{0.5, 1.0, 2.0\}$.
\textbf{Multi-architecture sweep}: CCI v2 fixed\_js and Model C are each
trained on DeBERTa-v3-base (184M), DeBERTa-v3-large (440M), and RoBERTa-large
(355M, different family), 6 seeds per cell (5 for one cell where one seed
was dropped). \textbf{Prompted LLMs (D-track)}: Mistral-7B-Instruct,
Qwen-2.5-7B-Instruct, Llama-3.1-8B-Instruct under three prompting strategies
(zero-shot, 5-example few-shot, Guided CoT~\citep{patel2025evaluating}) with
greedy decoding.

\subsection{Statistical protocol}

We pre-specify two test families and apply Holm-Bonferroni step-down within
each at $\alpha{=}0.05$~\citep{dror2018hitchhiker}. \emph{Family 1 (main)}
has 9 comparators $\times$ 6 metrics = 54 tests; \emph{Family 2
(within-architecture transfer)} has 2 pairs $\times$ 3 CFR metrics = 6 tests.
Paired-bootstrap by seed ($n_{\text{resample}}{=}10{,}000$) is the primary
test; we additionally report the exact sign-test p-value as a small-$n$
backup, which floors at 0.0313 for $n{=}6$ and 0.0625 for $n{=}5$.

\section{In-Distribution Results}
\label{sec:results-id}

\subsection{CCI v2 reduces CFR\textsubscript{hard} across three architectures (RQ1, RQ3)}

\begin{table}[h]
\centering
\small
\begin{tabular}{l r r r r r}
\toprule
Method & Params & F\textsubscript{1}\textsuperscript{+} $\uparrow$ & ECE $\downarrow$ & CFR\textsubscript{hard} $\downarrow$ & FNED $\downarrow$ \\
\midrule
Model C (DeBERTa-v3-base)            & 184M & $0.660 \pm 0.019$ & $0.077 \pm 0.034$ & $0.036 \pm 0.008$ & $0.797 \pm 0.159$ \\
Davani CLP $\lambda{=}0.5$           & 184M & $0.659 \pm 0.012$ & $0.066 \pm 0.018$ & $0.040 \pm 0.005$ & $0.786 \pm 0.036$ \\
\textbf{CCI v2 fixed\_js (pivot)}    & 184M & $\mathbf{0.677 \pm 0.008}$ & $0.098 \pm 0.053$ & $\mathbf{0.018 \pm 0.003}$ & $0.793 \pm 0.180$ \\
CCI v2 fixed\_js (DeBERTa-large)     & 440M & $\mathbf{0.719 \pm 0.012}$ & $0.099 \pm 0.063$ & $0.019 \pm 0.009$ & $0.647 \pm 0.058$ \\
CCI v2 fixed\_js (RoBERTa-large)     & 355M & $0.694 \pm 0.005$ & $0.077 \pm 0.052$ & $\mathbf{0.013 \pm 0.006}$ & $0.759 \pm 0.115$ \\
\bottomrule
\end{tabular}
\caption{In-distribution headline results, mean $\pm$ std over 6 seeds (5 for the
DeBERTa-large CCI cell after one seed was dropped). CCI v2 reduces
CFR\textsubscript{hard} by 49--73~\% versus same-architecture Model C
(DeBERTa-base: $-49\%$; DeBERTa-large: $-55\%$; RoBERTa-large: $-73\%$) with
F\textsubscript{1}\textsuperscript{+} preserved or improved.}
\label{tab:main}
\end{table}

All 12 base-architecture CFR comparisons (4 base comparators $\times$ 3 CFR
metrics) reject Holm-Bonferroni in the pivot's favour with $6/6$ seed
direction agreement. Within-architecture transfer tests (CCI v2 large vs
Model C large on the same backbone) reject 6/6 under Holm-Bonferroni; the
RoBERTa-large cell also passes the strict sign-test floor.

\subsection{A structural lemma: scalar TS cannot improve CFR\textsubscript{hard} (RQ2)}
\label{sec:lemma}

\begin{lemma}[Scalar TS preserves CFR\textsubscript{hard} at threshold $\tfrac{1}{2}$]
\label{lem:ts-cfr}
Let $f:\mathcal{X}\to\mathbb{R}^2$ be any binary classifier with positive-class
logit difference $z(x)$ and let $\tilde z(x){=}z(x)/T$ for any $T{>}0$. For any
counterfactual pair $(x, x')$,
\[
\mathbb{1}[\sigma(z(x))\geq\tfrac{1}{2}]\neq\mathbb{1}[\sigma(z(x'))\geq\tfrac{1}{2}]
\iff
\mathbb{1}[\sigma(\tilde z(x))\geq\tfrac{1}{2}]\neq\mathbb{1}[\sigma(\tilde z(x'))\geq\tfrac{1}{2}].
\]
\end{lemma}

\begin{proof}
$\sigma(\cdot)\geq\tfrac{1}{2}\iff(\cdot)\geq 0$, and dividing by $T{>}0$ preserves sign.
\end{proof}

The Lemma's content is one line, but its operational value is a regression test:
we verify across all six Model C seeds that post-hoc TS produces CFR\textsubscript{hard}
\emph{bit-exactly equal} to the raw classifier (fitted $T\in[0.566, 1.400]$,
mean 1.031). CFR\textsubscript{hard} reduction therefore requires either a
training-time intervention (CCI v2) or a non-monotone post-hoc method.
The latter has its own cost (\S\ref{sec:tension}).

\subsection{Multicalibration vs counterfactual invariance: a tension (RQ1, RQ2)}
\label{sec:tension}

We fit a per-input-group multicalibration patcher~\citep{hebertjohnson2018multicalibration}
on Model C's dev-set predictions and apply it at inference. The result:

\begin{table}[h]
\centering
\small
\begin{tabular}{l c c c c}
\toprule
& F\textsubscript{1}\textsuperscript{+} & ECE & CFR\textsubscript{hard} & CFR\textsubscript{soft, L\textsubscript{1}} \\
\midrule
Model C (raw)                     & 0.660 & 0.077 & 0.036 & 0.032 \\
Model C + multicalibration        & 0.652 & \textbf{0.035} & 0.033 & \textbf{0.063} \\
$\Delta$ (\%)                      & $-1\%$ & $-54\%$ & $-8\%$ & $\mathbf{+96\%}$ \\
\bottomrule
\end{tabular}
\caption{Multicalibration cuts ECE by 54~\% but \emph{inflates}
the L\textsubscript{1}-soft counterfactual flip rate by 96~\%. The mechanism:
multicalibration patches predictions within group-cells of the input space;
counterfactual pairs $(x, x')$ live in \emph{different} cells by construction,
so the patcher can shift them apart. To our reading this is the first
quantitative report of this trade-off on a text classifier; the general
calibration-vs-equal-error fairness incompatibility is established
by~\citet{pleiss2017fairness}.}
\label{tab:tension}
\end{table}

\subsection{Prompted LLMs --- a clean negative result (RQ5)}
\label{sec:results-llm}

\begin{table}[h]
\centering
\small
\begin{tabular}{l c c c}
\toprule
Model & D1 zero-shot & D2 few-shot (5 ex.) & D3 Guided CoT \\
\midrule
Llama-3.1-8B-Instruct  & 0.432 & 0.399 & \textbf{0.445} \\
Mistral-7B-Instruct    & 0.296 & 0.386 & 0.387 \\
Qwen-2.5-7B-Instruct   & 0.335 & 0.366 & 0.301$^{\dagger}$ \\
\midrule
\textit{Model C (fine-tuned baseline)} & \multicolumn{3}{c}{\textit{0.660}} \\
\textit{CCI v2 fixed\_js (pivot)}      & \multicolumn{3}{c}{\textit{0.677}} \\
\bottomrule
\end{tabular}
\caption{Positive-class F\textsubscript{1} for three prompted LLMs under three
prompting strategies, on the same 1{,}060-tweet test set. Greedy decoding
(temperature$=0$). The best LLM--prompt combination (Llama + Guided CoT)
reaches 0.445 --- 22 F\textsubscript{1}\textsuperscript{+} points below
the 184M-parameter fine-tuned baseline. $^{\dagger}$Qwen + CoT predicts
``antisemitic'' for 1055/1060 examples (recall 1.0, precision 0.18); the
CoT scaffold over-primes its specificity.}
\label{tab:llm}
\end{table}

Guided CoT helps Llama (${+}0.013$) and Mistral (${+}0.091$) over zero-shot
but degrades Qwen (${-}0.034$). For Llama, the CoT scaffolding triggers
20/1060 ($\approx 2\%$) safety-filter refusals on slur-containing posts;
the parser counts these as invalid and excludes them from the metric.

\section{Out-of-Distribution Results: Cross-Dataset Transfer (RQ4)}
\label{sec:results-ood}

We zero-shot-evaluate all six (model, architecture) cells on the Jewish-target
subsets of HateXplain~\citep{mathew2021hatexplain} (2{,}543 examples, 68~\%
positive base rate) and ToxiGen~\citep{hartvigsen2022toxigen} (684 examples,
43~\% positive).

\begin{table}[h]
\centering
\small
\begin{tabular}{l l c c c c c c}
\toprule
& & \multicolumn{3}{c}{HateXplain Jewish} & \multicolumn{3}{c}{ToxiGen Jewish} \\
\cmidrule(lr){3-5}\cmidrule(lr){6-8}
Method & Arch. & F\textsubscript{1}\textsuperscript{+}@0.5 & F\textsubscript{1}\textsuperscript{+}@opt & PR-AUC & F\textsubscript{1}\textsuperscript{+}@0.5 & F\textsubscript{1}\textsuperscript{+}@opt & PR-AUC \\
\midrule
Model C  & base  & $\mathbf{0.718}\pm0.039$ & $0.799\pm0.015$           & $\mathbf{0.803}\pm0.027$ & $\mathbf{0.521}\pm0.057$ & $0.668\pm0.024$           & $0.752\pm0.024$ \\
CCI v2   & base  & $0.683\pm0.082$           & $\mathbf{0.811}\pm0.001$ & $0.796\pm0.023$           & $0.480\pm0.074$           & $\mathbf{0.675}\pm0.023$ & $0.746\pm0.022$ \\
Model C  & DL    & $0.762\pm0.023$           & $0.809\pm0.007$           & $0.812\pm0.012$           & $\mathbf{0.588}\pm0.037$ & $\mathbf{0.713}\pm0.014$ & $0.798\pm0.015$ \\
CCI v2   & DL    & $\mathbf{0.775}\pm0.016$ & $\mathbf{0.811}\pm0.008$ & $\mathbf{0.832}\pm0.012$ & $0.558\pm0.034$           & $0.702\pm0.024$           & $\mathbf{0.801}\pm0.011$ \\
Model C  & RL    & $0.661\pm0.042$           & $0.798\pm0.013$           & $0.762\pm0.022$           & $\mathbf{0.587}\pm0.038$ & $0.677\pm0.019$           & $0.756\pm0.022$ \\
CCI v2   & RL    & $\mathbf{0.672}\pm0.033$ & $0.798\pm0.022$           & $\mathbf{0.782}\pm0.013$ & $0.540\pm0.064$           & $\mathbf{0.691}\pm0.030$ & $\mathbf{0.763}\pm0.022$ \\
\bottomrule
\end{tabular}
\caption{Cross-dataset transfer. DL = DeBERTa-v3-large; RL = RoBERTa-large.
F\textsubscript{1}\textsuperscript{+}@0.5 uses the default decision threshold;
F\textsubscript{1}\textsuperscript{+}@opt selects the threshold that maximises
F\textsubscript{1}\textsuperscript{+} on the same OOD test set (an
upper-bound diagnostic, not a deployable rule). Mean $\pm$ std over 5 or 6 seeds.}
\label{tab:xds}
\end{table}

\paragraph{What this shows.} At the default threshold $t{=}0.5$, CCI v2 loses
F\textsubscript{1}\textsuperscript{+} on 4 of 6 (architecture, dataset) cells
(by $-0.030$ to $-0.047$), wins by small margins on the two HateXplain large
cells ($+0.011$ to $+0.013$), and CCI v2 wins or ties PR-AUC on every cell
(per-cell $\Delta$ range $[-0.007, +0.021]$, CCI v2 winning on 4 of 6 iso-architecture cells).
The ranking quality is preserved cross-dataset; the operating point is not.

\paragraph{Threshold-sweep diagnostic.} For each cell we sweep the decision
threshold over $[0.05, 0.95]$ in steps of $0.01$ and report the
F\textsubscript{1}\textsuperscript{+}-maximising threshold $t^*$ per seed.
Two patterns emerge. First, \textbf{CCI v2's optimal $t^*$ is systematically
low}: median $t^* \in [0.05, 0.15]$ across all six CCI v2 cells, versus
Model C's $[0.05, 0.18]$ --- both are well below the default $0.5$. Second,
\textbf{CCI v2's threshold gain $\Delta=$F\textsubscript{1}\textsuperscript{+}@opt
$-$ F\textsubscript{1}\textsuperscript{+}@0.5 is larger at base scale}: the
base CCI v2 cell averages $+0.162$ of threshold gain on the two OOD datasets,
versus Model C base's $+0.115$. CCI v2 has more performance locked behind a
mis-calibrated default cutoff; at large scale this gap narrows
(${+}0.090$ vs ${+}0.086$ on DeBERTa-large).

\paragraph{Iso-architecture comparison at F\textsubscript{1}\textsuperscript{+}@opt.}
Once each model uses its own PR-curve-optimal threshold:
\begin{itemize}[leftmargin=*,itemsep=1pt]
    \item HateXplain converges across all six cells to F\textsubscript{1}\textsuperscript{+}$\in[0.798, 0.811]$
    --- a band of $0.013$. CCI v2 marginally wins on DeBERTa-base ($+0.012$),
    is tied within $0.002$ on DeBERTa-large, and is identical on RoBERTa-large.
    \item ToxiGen splits cleanly: CCI v2 wins on DeBERTa-base ($+0.007$) and
    RoBERTa-large ($+0.014$), loses by $-0.011$ on DeBERTa-large (within seed noise).
    Out of six iso-architecture comparisons, CCI v2 wins or ties five.
\end{itemize}

\paragraph{Cross-dataset bottom line.} The apparent F\textsubscript{1}\textsuperscript{+}
deficit of CCI v2 at $t{=}0.5$ \emph{is} a threshold-shift artefact, full stop.
CCI v2's training compresses scores toward the swap-pair mean, leaving the
model at a more conservative operating point cross-dataset; once each model
uses its own PR-curve-optimal threshold, CCI v2 wins or ties Model C on
five of six iso-architecture cells and PR-AUC is preserved or improved
on every cell (per-cell $\Delta$ range $[-0.007, +0.021]$).  The artefact is the predicted behaviour of a
Jensen-Shannon pair-invariance regulariser
(mechanism in \S\ref{sec:results-id}) and is consistent with the
calibration-vs-fairness incompatibility
of~\citet{pleiss2017fairness}.

\paragraph{The gap is calibration-driven, not representation-driven.}
F\textsubscript{1}\textsuperscript{+}@opt is oracle-tuned on the OOD set
itself and thus an upper-bound diagnostic. To isolate \emph{where} the
gap actually lives, we ran a deployable middle column: tune the threshold
on the GoldStandard2024 source test split (no OOD labels), apply it
unchanged to HateXplain and ToxiGen. Result: the source-tuned threshold
clusters in $[0.44, 0.62]$ across all twelve cells --- essentially the
default $0.5$ --- while the OOD-optimal threshold clusters in
$[0.05, 0.18]$. The $\approx 0.4$-unit gap is structural: the source has
$17.3\%$ positive base rate (well-calibrated at $\sim 0.5$), the OOD
sets have $43\%$ and $68\%$ (well-calibrated far below). A representation
deficit would manifest at \emph{any} threshold; the fact that the gap
disappears exactly once threshold matches OOD prevalence pinpoints it
as a prevalence-driven calibration shift, orthogonal to the ranking
quality the regulariser was designed to deliver
(supplementary \texttt{threshold\_sweep\_cross\_dataset.csv}, $n{=}68$
cells). Independently, $46.6\%$ of ToxiGen Jewish positives contain no
swappable identity token (vs $11.7\%$ of HateXplain Jewish positives), so
the swap-invariance training mechanism cannot operate on nearly half of
ToxiGen's positives; that this lexical limitation does not produce a
larger ToxiGen gap suggests the regulariser's invariance signal
generalises beyond the swap vocabulary.

\section{Discussion}

\paragraph{Three axes are genuinely orthogonal.} The Lemma rules out scalar
Temperature Scaling as a CFR\textsubscript{hard} fix; the multicalibration
finding shows that fixing subgroup-conditional ECE actively damages
soft counterfactual stability ($+96\%$). Conversely, CCI v2 reduces
CFR\textsubscript{hard} but does not improve global ECE and shows mild
worst-group ECE penalties. The three axes accuracy / calibration /
counterfactual stability are not aliases for each other, and a model
selection procedure that targets one will not automatically deliver on
the others.

\paragraph{Ranking transfers; the operating point shifts with prevalence.}
The cross-dataset evaluation pinpoints what training-time counterfactual
invariance buys out-of-distribution and what it does not. The regulariser
delivers \emph{rank-preserving invariance}: PR-AUC and
F\textsubscript{1}\textsuperscript{+}@opt are preserved or marginally
improved on every iso-architecture cell, with CCI v2 winning or tying
on 5 of 6 cells at $t^*$. It does \emph{not} deliver
decision-preserving invariance at a fixed default cutoff: the
Jensen-Shannon pair-invariance term pulls predictions on a swap pair
$(x, x')$ toward their geometric mean, compressing scores toward the
distribution centre. The right \emph{order} of antisemitic vs
non-antisemitic content is learned; the cut-off to convert that order
into decisions has to track the deployment distribution's positive
base rate (GoldStandard2024 $17\%$, ToxiGen $43\%$, HateXplain $68\%$).
This is exactly the scope \citet{veitch2021counterfactual} describe:
counterfactual invariance is defined against a specified intervention
set, with no claim about decisions at a fixed cutoff outside that set.
Our threshold-sweep + source-tuned check make the partition explicit:
the regulariser delivers rank-preserving invariance; deployment then
needs a cheap OOD-aware threshold pick.

\paragraph{Prompted LLMs underperform fine-tuned by 22 points.} This replicates
\citet{patel2025evaluating}'s general finding on the Jewish-target subset
and across three independent LLM families. Guided CoT helps Llama and
Mistral, destroys Qwen --- the kind of model-specific prompt brittleness
the deployability literature has been documenting~\citep{turpin2023language}.

\paragraph{Limitations.} (i) Single dataset (GoldStandard2024) for training;
(ii) single language (English); (iii) single identity target (antisemitism);
(iv) the swap engine is rule-based; an LLM-likelihood-based filter would
likely improve coverage on borderline pairs.

\section{Scope and Novelty Self-Assessment}

Honest claim-by-claim ranking of the work's novelty, on a 1--5 scale.

\begin{itemize}[leftmargin=*,itemsep=2pt]
\item \textbf{Multicalibration vs counterfactual-invariance tension (\S\ref{sec:tension}).}
\textbf{Novelty: 5/5.} To our reading this is the first quantitative report
of a tension between post-hoc multicalibration and counterfactual
pair-invariance on a text classifier. The $+96\%$ CFR\textsubscript{soft,L1}
regression is mechanistically traced to the partition structure of
multicalibration patchers. The general calibration-vs-fairness incompatibility
(\citet{pleiss2017fairness}) is established literature; the empirical
specialisation to text counterfactual pairs is new.

\item \textbf{Multi-architecture cross-family transfer of CCI.}
\textbf{Novelty: 4/5.} Demonstrating that the regulariser transfers
\emph{both} same-family up in scale (DeBERTa-base$\to$large, $+1.4\times$
parameters) \emph{and} cross-family (DeBERTa-base$\to$RoBERTa-large,
different tokenizer) is uncommon in the counterfactual-fairness
literature. The cross-family RoBERTa-large reduction
($-73\%$ CFR\textsubscript{hard}) is the strongest single empirical result.

\item \textbf{Two-family Holm-Bonferroni significance suite with paired
bootstrap and exact-sign-test floor.}
\textbf{Novelty: 4/5} for a course mini-project (most don't bother),
\textbf{2/5} as methods contribution in absolute terms (the design follows
\citet{dror2018hitchhiker}).

\item \textbf{Honest FNED Pareto audit + transparent pivot change.}
\textbf{Novelty: 3/5.} The original pre-registered pivot
(CCI v2 learned-$T$) was switched to fixed-$T$ after a within-cell
ablation showed the learned-$T$ cell cost FNED without buying
CFR\textsubscript{hard}. Most papers wouldn't admit this.

\item \textbf{Cross-dataset transfer analysis at three operating points.}
\textbf{Novelty: 4/5.} The cross-dataset evaluation reports
F\textsubscript{1}\textsuperscript{+} at three thresholds per
(architecture, dataset) cell: the default $t{=}0.5$ (deployed today), the
source-tuned $t^{\mathrm{src}}$ (no OOD labels), and the OOD-optimal
$t^*$ (upper bound). The three columns disambiguate (i) the artefact at
the default cutoff, (ii) the prevalence-driven structural component
($t^{\mathrm{src}}\!\approx\!0.5$ vs $t^*\!\approx\!0.1$), and (iii) the
preserved ranking ceiling (CCI v2 wins or ties on 5/6 iso-architecture
cells at $t^*$, PR-AUC $\Delta \in [-0.007, +0.021]$ winning on 4/6).
The three-operating-point
disambiguation is, to our reading, novel in pair-invariance cross-dataset
evaluation and gives a clean mechanistic diagnosis: the regulariser
preserves ranking and the gap at the default cutoff is a prevalence-shift
calibration issue, not a representation deficit.

\item \textbf{CCI v2 loss formulation (JS in probability space + per-group $T$).}
\textbf{Novelty: 3/5.} Incremental over Counterfactual Logit
Pairing~\citep{garg2019counterfactual,davani2021improving} which uses
$L_p$ in logit space. The JS choice yields different gradient dynamics
near the decision boundary; the per-group $T$ does not improve
CFR\textsubscript{hard} in our ablation but is retained as a comparator.

\item \textbf{Lemma~\ref{lem:ts-cfr} (scalar TS preserves
CFR\textsubscript{hard} at threshold $1/2$).}
\textbf{Novelty: 2/5} for the math (the proof is one line: monotone
$\sigma$ preserves sign across the threshold), \textbf{4/5} for the
operational regression-test verification (bit-exact preservation on six
Model C seeds with fitted $T\in[0.566, 1.400]$).

\item \textbf{Prompted-LLM negative result (best LLM 22 points below
fine-tuned across three LLM families).}
\textbf{Novelty: 2/5.} Replication of \citet{patel2025evaluating}'s
general finding on the Jewish-target subset; small extension via
Qwen-CoT brittleness observation.

\item \textbf{Bimodal ECE under focal + F\textsubscript{1}-based early
stopping (supporting observation).}
\textbf{Novelty: 1/5.} Already characterised by \citet{mukhoti2020calibrating}
and \citet{tao2023pcs}; we document it on this dataset.
\end{itemize}

\paragraph{Aggregate scope.} \textbf{Most novel}: the multicalibration vs
counterfactual-invariance tension (\S\ref{sec:tension}). \textbf{Most useful}:
the three-axis evaluation framework as a methodological habit for hate-speech
classifier evaluation, plus the three-operating-point (default / source-tuned
/ OOD-optimal) cross-dataset diagnostic that separates ranking transfer
from prevalence-driven calibration shift. \textbf{Most honest}:
the FNED Pareto audit and the cross-dataset three-operating-point analysis
that re-interprets our earlier ``negative result'' as a calibration shift,
not a representation deficit.

\section{Future Work}

\begin{enumerate}[leftmargin=*,itemsep=2pt]
\item \textbf{Held-out OOD dev split for deployable threshold calibration.}
Our source-tuned column shows that GoldStandard2024-side calibration alone
cannot close the OOD gap (prevalence shift dominates). The natural next step
is a small held-out OOD dev split per target distribution; the question is
how many labelled OOD examples are needed for the calibrated threshold to
generalise --- a sample-complexity ablation that fits cleanly in a follow-up.
\item \textbf{Composing CCI v2 with random identity-token masking.} A small
sweep replacing the pure-CCI cell with CCI~v2 $+$ Bernoulli masking
augmentation. Prediction: random masking provides omnidirectional smoothness
on examples without an explicit identity token (the $46.6\%$ ToxiGen
no-token positives) while CCI v2's directional pair-invariance handles the
within-vocabulary cases.
\item \textbf{Mechanistic derivation of the score-compression effect.}
The Jensen-Shannon pair-invariance gradient
$\partial\mathrm{JS}/\partial z \propto p(1{-}p)\,[\mathrm{logit}(p)-\mathrm{logit}(m)]$
formally predicts the contraction-toward-mean behaviour we observe empirically
(low $t^*$ for CCI v2). Writing this out and verifying empirically would tie
the threshold-sweep finding to a closed-form mechanism.
\item \textbf{LLM-likelihood-based symmetry filter.} The plumbing exists;
plugging it in would expand the swap manifest's coverage.
\item \textbf{Cross-language and cross-identity transfer.} The swap engine
and the symmetric vocabulary are the only identity-specific components; the
method is, in principle, identity-agnostic.
\end{enumerate}

\section{Conclusion}

Training-time counterfactual invariance is a real,
mechanistically-interpretable axis of regularisation for hate-speech
classifiers, orthogonal to accuracy and calibration.
\textbf{In distribution}, CCI v2 cuts the hard counterfactual flip rate
by 49--73~\% across three architectures, with
F\textsubscript{1}\textsuperscript{+} preserved or improved, and we
identify a previously-undocumented tension between multicalibration and
counterfactual pair-invariance.
\textbf{Out of distribution}, the gain transfers cleanly as preserved
ranking quality (PR-AUC range $\Delta \in [-0.007, +0.021]$ across six
iso-architecture cells, CCI v2 winning on 4 of 6) at a more conservative operating point
($t^*\!\approx\!0.05$--$0.15$ vs the default $0.5$); at each model's
PR-curve-optimal threshold, CCI v2 wins or ties on five of six
iso-architecture cross-dataset cells. A source-tuned threshold check
($n{=}68$ cells) locates the residual gap at default cutoff in a
$\approx 0.4$-unit prevalence shift on the optimal threshold rather
than in the representation. This separates two distinct claims about
generalisation that the literature often conflates:
\emph{ranking generalisation} (which the regulariser delivers
out-of-the-box) and \emph{decision generalisation at a fixed default
cutoff} (which requires a small OOD-aware threshold calibration step).

\small
\bibliographystyle{plainnat}
\bibliography{references}

\end{document}
